With the initial model outlined, we are now in a position to review it for its limitations and potential improvements. The model is a simplified representation of a complex system, and as such it has several points that could be improved upon. However, it serves as a useful conversation starter, with a whole slew of questions being raised in the process of creating it.

We will start by reviewing the model, discussing the potential faults and improvements, followed by the implications and questions raised by it. We will then conclude with a few cases that could be interesting for the model to be applied to, which would in turn change how the model is structured.

\subsection{Review of the Model}
\label{ReviewOfModel}

For the review of our model, we have many topics to cover, as there are numerous potential improvements that could be made to it. The model is a very simplified representation of a complex system, and as such it has several limitations that need to be addressed. We will be looking first at how time plays a significant role in the interactions between the agents and the system, and how the model does not currently account for this, then look more critically at the states that we have defined, and how they might not be entirely representative of the system. 

We will look at the lack of temporary memory allocation representation in the model, such as RAM. Then we will look at the question of privilege level, and the lack of explicit authorisation constraints in the model. We then examine the symmetry of the agents and how this assumption may affect the model. Lastly, we will discuss the potential for a threat vector analysis to be applied to the model, and why it was ultimately not used in this case.

\subsubsection{Time Series}
\label{TimeSeries}

The model we have composed is essentially a snapshot of a single moment in time. It does not take into account the temporal aspect of the problem, which is a significant limitation in this study, given that a great deal of the interactions between the agents and the system are time-dependent. For example, the time it takes for an agent to think of its next move or the time it takes for a command to execute can have a significant impact on the strategies an agent might employ.

Seeing as inodes are indexed as well, we could also see that it might be possible to track the history of file access and modification quite a lot faster than many other actions on the system. Since every file operation updates inode timestamps, a rival agent could use commands like \texttt{find} with a \texttt{-newer} flag, or \texttt{stat} on individual files, to rapidly reconstruct a timeline of recent changes. This raises the question of whether some of the agent strategies are rendered obsolete by this fact. If a rival agent can string together a short sequence of commands to enumerate every file modified in the last few seconds, newly created backups and dummy files would be immediately visible by virtue of their timestamps, potentially making those strategies far less effective than the analysis suggests. 

In order for the model to work with time series, we would need to have continuous snapshots of the system's state to see how it changes over time, which might quickly make the model difficult to interpret and analyse, defeating the purpose. However, it is possible to see how a system changes over time, and by selecting a choice amount of snapshots, we can still get a good idea of how the system evolves without overwhelming ourselves with data. 

\subsubsection{States}
\label{StatesDiscussion}

The model we have created is based on a set of states that represent the condition of a set of files in the system. This of course is subject to a few flaws, as it requires an abstract representation of the system, which might not be entirely representative of the actual system. For example, dead agents are represented equally in the model, but in reality there might be nuances to how an agent gets killed, such as actual formatting, overwriting, or just deleting the file. These different methods of killing an agent might have different implications for the rival agent, as these all have vastly different time implications, potentially making some strategies more viable than others.

More broadly, the binary alive/dead representation is a significant simplification. In practice, an agent could exist in a range of intermediate states: partially corrupted, where some but not all of its files have been modified; degraded, where its main process is still running but one or more watchdogs have been eliminated; or locked out, where the agent is alive but cannot access the filesystem due to permission changes. Each of these intermediate states would carry different strategic implications for both the agent and its rival, and collapsing them all into a single "dead" state means the model loses information that could be relevant to understanding how and why one strategy outperforms another.

\subsubsection{Multiple Agents}
\label{MultipleAgents}

The model currently only accounts for two agents, which in itself is a significant change in how agents are viewed in the context of a system. In theory, there could be multiple agents running on the system simultaneously, which would make the interactions between them much more complex. With three or more agents, the notion of purely adversarial competition breaks down: it becomes necessary to consider the possibility of temporary alliances, where two agents cooperate to eliminate a third before potentially turning on each other. This introduces a game-theoretic dimension well beyond the zero-sum framing of the current two-agent model, and raises questions about how an agent would decide when to cooperate and when to defect.

In addition to this, multiple agents would make the user profile system more strategically relevant. Linux provides built-in mechanisms for switching between user profiles, and commands such as \texttt{write} and \texttt{wall} allow messages to be sent between logged-in users. These could serve as a primitive communication channel between agents operating under different user accounts, opening up coordination strategies that do not leave traces in the filesystem at all. This would also require the model to account for multiple privilege domains running in parallel, further complicating the privilege level discussion above.

\subsubsection{Lack of Temporary Memory Allocation}
\label{TemporaryMemory}

In addition to this, the process space used by the agents is represented in the model as states, but the model does not account for temporary memory allocation in RAM. This is a significant gap, as RAM and disk storage have fundamentally different characteristics that would influence agent strategy. RAM is orders of magnitude faster to read and write than disk, but it is volatile: its contents are lost when a process terminates or the system is rebooted. An agent that stores working data or intermediate results in memory rather than on disk would be much harder for a rival to detect through filesystem monitoring, since nothing is written to a path that tools like \texttt{ls} or \texttt{find} would surface. At the same time, memory usage is visible through tools like \texttt{top} and \texttt{htop}, meaning that an agent making heavy use of RAM would still leave a detectable footprint, just in a different domain.

The \texttt{/proc} pseudo-filesystem is related but distinct. Rather than being a storage medium, \texttt{/proc} is a kernel-provided interface that exposes live information about running processes, including their memory maps, open file descriptors, and environment variables. An agent could read from \texttt{/proc} to gather intelligence about a rival process without writing anything to disk, making it an inherently stealthy surveillance channel. How agents might exploit this is discussed further at the end of this section.

\subsubsection{Privilege Level}
\label{PrivilegeLevel}

One variable that the current model leaves unresolved is the privilege level at which the agents operate. The analysis assumes that both agents can freely read and write files, spawn and terminate processes, and access system information, but whether they are running as an unprivileged user or as root has a significant impact on what is actually possible. An unprivileged agent cannot terminate processes owned by another user, cannot write to system directories, and cannot modify file permissions it does not own. A root-level agent, on the other hand, faces almost none of these restrictions, which substantially widens the set of viable strategies.

This distinction matters for the model in a concrete way. Strategies such as killing the rival agent's processes or overwriting its files presuppose a level of access that may not be granted in a realistic deployment scenario. Future iterations of the model should specify the privilege context explicitly, and may benefit from modelling the two cases separately, as the set of available actions and therefore the dominant strategies are likely to differ considerably between them.

\subsubsection{Symmetry of the Agents}
\label{AgentSymmetry}

The model as constructed treats the two agents as symmetric: equal in resources, capabilities, and knowledge of the system. This is a useful simplifying assumption for building an initial model, but it does not reflect the more likely real-world scenario, where an attacker and a defender are rarely evenly matched. An attacker agent may have been designed with specific knowledge of the defending system, may have more computational resources available, or may simply have the advantage of striking first while the defender is still establishing its baseline state map.

Conversely, a defender may have the advantage of being able to set up monitoring and integrity checks before the attacker is introduced, giving it an informational head start. Asymmetry in either direction changes the dynamics of the simulation considerably, and acknowledging this as a limitation of the current model is important for contextualising its conclusions. A more advanced model could introduce configurable asymmetry as a parameter, allowing the simulation to explore a range of scenarios rather than just the balanced case.

\subsubsection{The Arms Race Dynamic}
\label{ArmsRaceDynamic}

A natural consequence of the strategies identified in the analysis is the emergence of an arms race dynamic between the two agents. As each agent develops defensive measures, the rival is incentivised to develop more sophisticated means of circumventing them, which in turn prompts further defensive refinement. This is not a new phenomenon in cybersecurity, but it is worth considering how it manifests specifically within the constraints of this model.

The key driver of the arms race here is resource cost. Every defensive measure an agent adds, whether watchdogs, backups, or dummy files, consumes CPU time, memory, and disk I/O that could otherwise be used for offensive actions or for faster state map updates. The rival agent, aware of this overhead, can respond by increasing the rate or volume of its own actions to overwhelm the defender's monitoring capacity before it can respond. The defender must then either accept a degraded response time or invest more resources in monitoring, which again raises its own overhead and visibility.

The question this raises is whether the arms race converges to a stable equilibrium or whether it escalates indefinitely. In theory, the finite resources of the system impose a ceiling on how much either agent can escalate, meaning the race must eventually reach a point where neither agent can meaningfully improve its position without compromising its own survival. In practice, the first agent to exhaust its resource budget is likely to lose, which suggests that the optimal strategy may not be the most aggressive or the most defensive, but the most resource-efficient. This is a hypothesis that the simulation could be designed to test directly.

\subsubsection{Threat Vector Analysis}
\label{ThreatVectorAnalysis}

During the development of this model, the use of a formal threat vector analysis was considered as a means of systematically enumerating the attack surface between the two agents. In a threat vector analysis, each asset is mapped against a set of potential attack methods, producing a structured matrix of vectors that can then be prioritised by impact and likelihood. Applied to the AESM, this would have meant mapping each state category against the specific commands and mechanisms available to a rival agent.

The reason this approach was ultimately set aside is that the model, in its current form, is not yet tangible enough to support a rigorous analysis of this kind. A meaningful threat vector analysis requires a well-defined system with known components, established trust boundaries, and a concrete set of agent capabilities. The AESM is still a conceptual framework rather than an implemented simulation, and the agents within it are abstractions rather than programs with fixed capabilities. Applying a formal analysis at this stage would produce results that are more speculative than informative, and could give a false impression of precision. As the model matures and the simulation is implemented, a threat vector analysis would become a natural and valuable next step.

\subsection{Implications and Questions Raised}
\label{ImplicationsAndQuestions}

Now that we have reviewed some of our major concerns for the model, we can move on to the implications and questions raised by it. The model, while simplified, still raises a number of interesting questions about the nature of agent interactions, the limits of surveillance and countermeasures, and the potential for emergent behaviour.

\subsubsection{VFS Bypass}
\label{VFSBypass}

A question raised during the research of the filesystem and the way memory is written to it, is whether it is possible for an agent to bypass the VFS layer and write directly to a raw block device, which would allow it to store data outside the normal directory structure and avoid detection by the rival agent's filesystem monitoring. If such a write were possible, the data would not be associated with any inode or directory entry, meaning it would not appear in the output of \texttt{ls}, \texttt{find}, or any other VFS-level tool. The rival agent would have no way of knowing whether hidden data exists unless it manually scanned every block on the device, which would be an extremely time-consuming process.

In practice, writing directly to a block device requires root privileges and carries a significant risk of corrupting the filesystem by overwriting blocks that are already in use by the VFS. This makes it a high-risk strategy that would likely only be viable under specific conditions, and it falls outside the scope of the current simulation for that reason. However, as a theoretical capability it is worth acknowledging, as it illustrates that the state map model has a fundamental blind spot: it is only as complete as the VFS's view of the disk, and any data that circumvents that layer is invisible to it.

\subsubsection{Starting Scenarios}
\label{StartingScenarios}

The model as it stands is designed to simulate a scenario where two agents are introduced to a system at the same time, with no prior knowledge of each other or the system. This is a useful baseline scenario for testing and learning, but it raises the question of how the agents would perform under different initial conditions. Starting scenarios are significant because the early phase of the simulation, before either agent has established a complete state map or deployed any defences, is likely to be the most strategically consequential. Small differences in initial conditions could produce very different outcomes.

\textbf{Head Start for the Defender:} In this scenario, the defending agent is introduced to the system before the attacker and is given time to establish its baseline state map, deploy watchdog processes, and distribute backup files before the rival appears. This more closely mirrors a realistic cybersecurity context, where a system's defences are in place before an intrusion occurs. The question for the model is how much of a head start is sufficient to give the defender a meaningful advantage. If the defender can complete its setup before the attacker begins scanning, it may be able to operate in a much stronger position; if the attacker is introduced before the defender has finished, the incomplete defences may be worse than none at all, as they consume resources without providing full protection.

\textbf{Undetected Infiltration:} In this scenario, the attacker is introduced to the system silently, without the defender's knowledge, and operates for a period before the defender becomes aware of its presence. This represents the more dangerous real-world case, where an attacker has already established a foothold by the time defensive measures are activated. The defender would be starting from a compromised baseline, meaning its initial state map would already contain processes or files belonging to the attacker that it might misidentify as legitimate. This scenario tests a different capability: rather than preventing infiltration, the defender must detect and evict an agent that has already embedded itself in the system, which requires distinguishing between its own artefacts and the attacker's with no clean reference point.

\textbf{Resource Imbalance:} A third scenario worth considering is one where the two agents begin with unequal access to system resources, either because they are running under different user accounts with different resource limits, or because the system itself is under load from other processes. An agent operating under tight resource constraints would need to make different decisions about whether to prioritise monitoring, defence, or offence, as it cannot afford to do all three simultaneously. This scenario would stress-test the resource-efficiency hypothesis raised in the arms race discussion, as it would force at least one agent to operate below the resource floor at which a balanced strategy is viable.